\subsection{Zero-density obstruction to a height-independent heat comparison}
\label{subsec:heatgap}
The bounded-in-$x$ comparison in
Proposition~\ref{prop:heatsufficient} is a valid
sufficient hypothesis, but for the genuine
Riemann heat flow its hypothesis is false
already at the initial time $t=\tfrac12$.
The obstruction does not depend on a
hypothetical nonreal zero. It follows from
the \emph{known real zeros} of $H_{1/2}$
and their increasing density.

\begin{lemma}[Laguerre ratios at a simple real zero]
\label{lem:gapratio}
Let $G$ be a nonzero real entire function
of order at most one with only real zeros.
Suppose $a$ is a simple zero and $b\ne a$
is another real zero. Then, for the
generalized Laguerre coefficients defined
in Section~\ref{sec:transformattack},
\begin{equation}
\frac{L_2[G](a)}{L_1[G](a)}
=-\left(\log\left|
\frac{G(z)}{z-a}\right|\right)''\!\!(a)
\ \ge\ \frac{1}{(a-b)^2}.
\label{eq:gapratio}
\end{equation}
Here the logarithmic derivative is taken
in a real neighborhood of $a$.
\end{lemma}
\begin{proof}
Factor $G(z)=(z-a)g(z)$, with $g(a)\ne0$.
The first three derivatives at $a$ are
\[
G'(a)=g(a),\qquad
G''(a)=2g'(a),\qquad
G'''(a)=3g''(a).
\]
The coefficients of $y^2$ and $y^4$
in $G(a+iy)G(a-iy)$ are therefore
\begin{align*}
L_1[G](a)&=g(a)^2,\\
L_2[G](a)&=g'(a)^2-g(a)g''(a).
\end{align*}
Taking their ratio gives the equality
in \eqref{eq:gapratio}.
For the lower bound, use the real-zero
canonical product of an entire function
of order at most one:
\[
G(z)=e^{A+Bz}
\prod_\rho E_1(z/\rho),
\qquad E_1(w)=(1-w)e^w,
\]
with the usual adjustment for zeros
at the origin. On a real neighborhood
of a simple zero $a$, removing its
factor and differentiating the
logarithm twice gives
\[
-\bigl(\log|g|\bigr)''(a)
=\sum_{\rho\ne a}\frac{m_\rho}{(a-\rho)^2}.
\]
The series converges locally;
all terms are nonnegative, and the
summand for $b$ gives
\eqref{eq:gapratio}.
\end{proof}

\begin{theorem}[The proposed bounded heat
comparison fails for the actual theta family]
\label{thm:heatgapobstruction}
Let
\[
H_t(z)=\int_0^\infty
e^{tu^2}\Phi(u)\cos(zu)\,du,
\]
with $\Phi$ the genuine Riemann theta kernel.
Then the zeros of $H_{1/2}$ are all real
and simple, but
\begin{equation}
\boxed{
\sup_{\substack{x\in\mathbb R\\
L_{1,1/2}(x)>0}}
\frac{6L_{2,1/2}(x)}
{L_{1,1/2}(x)}=+\infty.
}
\label{eq:unboundedratio}
\end{equation}
Indeed, for arbitrarily large $T$, there
is a positive zero $a\in[T,2T]$ with
\begin{equation}
\frac{6L_{2,1/2}(a)}
{L_{1,1/2}(a)}
\ge\left(\frac{3}{8\pi^2}
+o(1)\right)(\log T)^2 .
\label{eq:loggaplower}
\end{equation}
In particular no finite value of
$c_1(0)$ can satisfy the
comparison \eqref{eq:comparisonmissing}
for all real $x$.
\end{theorem}
\begin{proof}
The de Bruijn--Newman theorem and the
strict upper bound $\Lambda<1/2$
ensure that $H_{1/2}$ has only real
simple zeros. Moreover, the known
zero-count asymptotics for each
fixed positive heat time give
\[
N_{1/2}(T)
=\frac{T}{4\pi}\log T+O(T).
\]
This is a classical input, not an
estimate established by the present
paper; compare Ki--Kim--Lee
\cite{KiKimLee} and the discussion of
fixed-time zero counts in
\cite{Polymath15,RodgersTao}.
Thus
\[
N_{1/2}(2T)-N_{1/2}(T)
=\left(\frac1{4\pi}+o(1)\right)
T\log T.
\]
Ordering the zeros in $[T,2T]$,
at least two adjacent positive zeros
$a<b$ have separation no greater than
\[
b-a\le\frac{T}{N_{1/2}(2T)
-N_{1/2}(T)-1}
=\frac{4\pi+o(1)}{\log T}.
\]
The transform $H_{1/2}$ has order
one and a real-zero canonical product.
Lemma~\ref{lem:gapratio} therefore gives
at the simple zero $a$,
\[
\frac{6L_{2,1/2}(a)}
{L_{1,1/2}(a)}
\ge\frac6{(a-b)^2}
\ge\left(\frac{3}{8\pi^2}
+o(1)\right)(\log T)^2 .
\]
Letting $T\to\infty$ proves
\eqref{eq:unboundedratio}.
\end{proof}

This result changes the logical
status of Proposition~\ref{prop:heatsufficient}:
its bounded comparison is not just
unverified but \emph{inapplicable to
the genuine theta heat flow}.
The earlier Gaussian test and the
double-zero example do not detect this
obstruction, which is present even
at a time where all actual theta-flow
zeros are known to be simple and real.
The theorem is an elementary
consequence of established zero-density
asymptotics and canonical products;
we make no novelty claim for those inputs.

\begin{proposition}[A corrected, spatially weighted
sufficient comparison]
\label{prop:weightedheatsufficient}
Suppose for each $n\ge1$ that a
nonnegative, locally bounded function
$c_n(s,x)$ on $[0,\tfrac12]\times\mathbb R$
satisfies
\begin{equation}
(n+1)(2n+1)U_{n+1}(s,x)
\le c_n(s,x)U_n(s,x)
\label{eq:weightedheatmissing}
\end{equation}
for every $s,x$, with $U_n$ as
defined above. Assume the usual
boundedness and regularity hypotheses
ensuring the whole-line parabolic
comparison principle. Then
$U_n(s,x)\ge0$ for every
$n\ge1$, $0\le s\le\tfrac12$,
and $x\in\mathbb R$.
Consequently all generalized
Laguerre inequalities hold for
$H_0$ and RH follows.
\end{proposition}
\begin{proof}
The hierarchy and
\eqref{eq:weightedheatmissing} give
\[
\partial_s U_n-\frac12\partial_x^2U_n
+c_n(s,x)U_n\ge0.
\]
Its initial value $U_n(0,x)$ is
nonnegative because $H_{1/2}$ has
only real zeros. The comparison
principle on the real line for
nonnegative locally bounded potentials,
under the stated decay and boundary
conditions, yields $U_n(s,x)\ge0$.
At $s=1/2$ the generalized
Laguerre--P\'olya criterion gives
the real-zero property of $H_0$.
\end{proof}

\paragraph{What a successful global estimate must look like.}
Theorem~\ref{thm:heatgapobstruction}
forces every admissible comparison
coefficient at $s=0$ to be
unbounded in the spatial variable,
and along a sequence to grow
at least as fast as $(\log |x|)^2$.
A possible sharper target would be
\[
0\le c_n(s,x)\le
C_n\bigl(1+\log^2(2+|x|)\bigr),
\qquad 0\le s\le\tfrac12,
\]
with constants $C_n$ independent
of $x$ but allowed to depend on $n$.
\emph{This estimate is a candidate,
not a theorem or verified conjecture.}
In particular it could fail at a
double-zero collision, where
$U_1=0<U_2$; the earlier
transversality calculation identifies
the exact mechanism.
The outstanding RH challenge is to
prove a valid spatially weighted
all-order comparison, or to obtain
a zero-exclusion argument that does
not require it.

\paragraph{Verification.}
The equality and inequality of
Lemma~\ref{lem:gapratio} are checked
at the local derivative level in
\path{XiKernel/HeatZeroGapObstruction.lean},
including a rational comparison gate.
The independent symbolic and
exact-rational replay in
\path{certificates/xi_heat_zero_gap_exact.py}
checks the local polynomial identity
and a sequence of shrinking-gap
quadratic models. The global
statement additionally uses the
published real-zero density
asymptotic for $H_{1/2}$
and the standard canonical product;
these substantial analytic inputs
are \emph{not} presently formalized
in Lean. The conclusion is a
rigorous negative test for one
proposed heat-flow proof step,
not an RH proof.
